\begin{cnabstract}
    三维光场显示技术作为下一代沉浸式视觉体验的核心载体，凭借裸眼立体感知优势在商业广告、医疗诊疗、影视娱乐等领域展现出广阔应用前景。光照作为构建立体真实感与空间可信度的关键因素，其编辑质量直接决定三维光场人像的显示效果与沉浸体验。然而，现有三维光场人像光照编辑方法存在采集设备复杂昂贵、多视角光照一致性不足、动态场景时序闪烁等技术痛点，难以满足高质量静态与动态光场内容的生成需求。
    
针对上述问题，本文围绕三维光场人脸光照编辑展开系统性研究，构建了从静态基础方法到动态优化框架的完整技术体系。主要研究工作与创新成果如下：

首先，提出一种关键信息驱动的三维光场静态人脸光照编辑方法。针对静态场景下多视图独立编辑导致的光照不一致与密集采集设备成本高昂问题，设计 “高光去除 - 光照编辑 - 密集视点生成” 三级框架：通过亮度阈值与漫反射分量双重检测机制精准定位并去除人像高光区域；采用双路径光照编辑策略，既支持从 HDR 图像提取球面调和（SH）系数实现环境光重光照，又能通过 BERT 模型解析自然语言指令，实现光照方向与强度的语义化调控；基于 NeRF 体渲染技术生成水平 ±30° 范围内 90 个视角的密集视差图，满足光场显示需求。实验结果表明，该方法在 FID（33.55）、PSNR（32.06）等指标上优于 NVPR、LRPI 等主流方法，单帧推理时间仅 25.92s，三维光场显示器上的立体感知得票率 89.19\%，光照一致性得票率 92.29\%。

其次，提出一种时间 - 空间双维联合优化的三维光场人脸视频动态光照编辑框架。针对静态方法扩展至视频序列时出现的时序闪烁与多视角光照失配问题，设计两大核心优化模块：时序一致性优化模块通过共享编码器提取稳定语义特征，结合光流对齐与光照 - 几何特征分离策略，施加 “特征 L1 + 像素残差 L2” 双重损失约束，有效抑制帧间光照闪烁；多视角一致性优化模块基于相机内外参实现三维几何对齐，通过三维表面采样将多视角光照约束转化为三维点光照观测问题，结合可见性、法线夹角、光照置信度三重权重的多因子聚合机制，实现跨视角光照的物理一致性表达。实验验证表明，该框架使时序稳定性指标 WE 降至 0.0004，多视角光照误差 LE 低至 0.7922，主观实验综合评分达 4.73，显著优于 PTI、SMFR 等对比方法。

本文通过理论创新与实验验证，有效解决了三维光场人脸光照编辑中的核心技术痛点，为高质量静态与动态光场内容生成提供了新的技术路径，相关成果可应用于裸眼 3D 显示、VR/AR、影视制作等领域，具有重要的理论意义与工程应用价值。

    \cnkeywords{三维光场显示；人脸光照编辑；球面调和函数；NeRF；时序一致性；多视角一致性}
\end{cnabstract}